\section{Worked problems: sizing and invalidating a cache}

\subsection{Llama 3 405B on eight H200s}
\textbf{Problem.} Llama~3~405B has 126 layers and 8 key/value heads of width
128 \cite{llama2024herd}. Compute its cache per token in BF16 and in FP8. Its
weights are stored in FP8 across eight H200s of 141~GB each; setting aside 4~GB
per GPU for activations and workspace, how many sequences of 131,072 tokens
fit?

\textbf{Solution.} Equation~\eqref{eq:kv-per-token}: $2 \times 126 \times 8
\times 128 \times 2 = 516{,}096$ bytes, 504~KiB per token in BF16, and 252~KiB
in FP8. A 131,072-token sequence therefore holds 66.1~GB of cache in BF16 or
33.0~GB in FP8. The eight GPUs have 1,128~GB; the FP8 weights take about
406~GB and the set-aside 32~GB, leaving about 690~GB. That is 10 sequences with
a BF16 cache or 20 with an FP8 cache (derived). Eight GPUs and a model that
reads 406~GB of weights per decode step serve ten or twenty long conversations
at a time: at this context the cache, not the weights, sets the batch, and
halving the cache's bytes doubles it. Spread over the eight GPUs by key/value
head, as tensor parallelism does (\Chref{model-parallelism}), each GPU holds
one of the eight heads of every layer.

\subsection{A batch of 64 on one H100}
\textbf{Problem.} Qwen3-8B in BF16 (16.4~GB of weights, 15.1~GB read per decode
step, 144~KiB of cache per token) serves 64 sequences on one 80~GB H100 with
4~GB set aside. At what context per sequence does a decode step read more cache
than weights? At what context does the batch no longer fit?

\textbf{Solution.} The cache read equals the weight read when
$64 \times S \times 147{,}456 = 15.1 \times 10^{9}$, at $S \approx 1{,}600$ tokens
per sequence. The memory for caches is $80 - 16.4 - 4 = 59.6$~GB, so the batch
fits while $64 \times S \times 147{,}456 \le 59.6 \times 10^{9}$, up to about
6,300 tokens per sequence (derived). Between the two, the batch fits but every
step reads more cache than weights; the machine is doing what batching was
supposed to prevent --- streaming bytes that serve one sequence each. Past
6,300 tokens the engine must shrink the batch, evict, preempt or move caches to
slower memory (\Chref{paged-kv}, \Chref{kv-storage-tier}).

\subsection{Editing one token of a long prompt}
\textbf{Problem.} A user edits the 1,000th token of a 4,096-token prompt that an
engine has already prefilled and cached. Which cache entries are still valid?
What does it cost on the M1 of \Chref{prefill-decode} to bring the cache up to
date, and what if the user had inserted a token instead of replacing one?

\textbf{Solution.} Entries for positions 0 to 998 are valid in every layer:
causal attention means nothing before the edit depends on it. The edited
token's own entries are invalid in every layer. The entries of later tokens are
valid in the first layer only: their layer-0 keys and values are projections of
their own embeddings, which did not change, but from layer 1 on their hidden
states include attention over the edited token. So the engine must recompute
the prefill for positions 999 to 4,095 --- 3,097 tokens, attending to up to 4,096
--- about four fifths of the full prefill's work (three quarters of the matrix
work, and 94\% of the attention, because the later tokens are the ones with the
most to attend to): roughly 14~s instead of 18.2~s with llama.cpp's default
attention path on the M1, or 12~s instead of 15.7~s with flash attention
(derived from the measured 4,096-token prefill). An
insertion is worse. Every later token moves one position, and because keys are
stored after their rotary rotation, even their first-layer keys change: every
entry from the insertion point on, in every layer, is invalid. This is why
prefix caches match exact token prefixes at exact positions and nothing more
(\Chref{prefix-caching}), and why an agent that appends to its context keeps its
cache while one that edits the middle of it pays for a new prefill.
